% =====================================================================
%  Software, data and legal status of this edition (back matter)
% =====================================================================
\chapter*{Software, data and legal status}\label{ch:status}
\phantomsection
\addcontentsline{toc}{chapter}{Software, data and legal status}
\markboth{Software, data and legal status}{Software, data and legal status}

Some of the content of this book will date faster than the rest: versions of software, online
services and the law. This chapter gathers it in one place, as it stood when this edition was
prepared (October 2026), so that the next edition, or a reader, can check it in a single list.

\section*{Software}

\begin{sloppypar}
Every listing was run with Python 3.12 and the packages of \cref{tab:versions}, and the outputs printed in the book are those of that run. The versions are pinned in the file
\code{code/requirements.txt}. To
reproduce the outputs, create a virtual environment with Python 3.12 and install the packages with
\begin{center}\code{pip install -r code/requirements.txt}\end{center} Other versions usually work, but numbers can change in
the last digits and some default parameters may differ.
\end{sloppypar}

\begin{table}[htbp]
\centering\small
\caption{Versions of the Python packages used to run the listings.}\label{tab:versions}
\begin{tabular}{@{}lll@{}}
\toprule
Package & Version & Used in\\
\midrule
numpy, scipy, pandas & 2.4.4, 1.17.1, 3.0.2 & all units\\
scikit-learn & 1.8.0 & all units\\
imbalanced-learn, matplotlib, joblib & 0.14.2, 3.10.8, 1.5.3 & all units\\
Flask, PyYAML & 3.1.3, 6.0.3 & Unit~5\\
pytest, hypothesis & 9.1.1, 6.169.0 & Unit~5\\
onnx, onnxruntime, skl2onnx & 1.23.2, 1.24.4, 1.20.0 & Unit~5\\
sklearn2pmml & 0.133.1 & Unit~5 (needs Java 17 or later)\\
shiny, starlette, uvicorn & 1.8.0, 1.7.0, 0.54.0 & Units~5 and~7\\
\bottomrule
\end{tabular}
\end{table}

Some examples need software outside Python: Java 17 or later for PMML (\cref{sec:exchange}), a
container engine such as Docker for packaging (\cref{sec:packaging}), and Git with a continuous
integration service, such as GitLab CI/CD, for the pipelines of \cref{sec:infrastructure}.

\section*{Data}

The German credit\index{German credit data} data are downloaded from a mirror pinned to a fixed commit,
and the listings check the SHA-256 hash of the file (\cref{lst:firstlook}). If the mirror disappears,
the same data are available from the UCI Machine Learning Repository \citep{hofmann1994statlog},
in a different file format. Most other examples generate synthetic data with fixed random seeds, and
the listings that download other data say where they come from in their code.

\section*{Services and models}

Online services change more often than libraries. Posit plans to move its older hosting service,
shinyapps.io, to Connect Cloud during 2027 (\cref{sec:shinyapps}), so the options for sharing an
application should be checked against the current documentation. Applications built on language
models depend on versions that providers update or retire (\cref{sec:llmmonitor}), and any result
obtained with a provider's model holds only for the version used, which must be recorded
(\cref{sec:llmsystems}).

\section*{The AI Act}

\Cref{sec:aiact} describes the AI Act\index{AI Act} as amended by the Digital Omnibus on AI
(Regulation (EU) 2026/1744), as checked on 12 October 2026. The dates that matter most for the uses
discussed in this book are these: the prohibitions and the duty to support AI literacy apply since
2 February 2025, and the obligations for general-purpose models since 2 August 2025; the obligations
for the high-risk systems of Annex~III, such as credit scoring, apply from 2 December 2027, and those
for high-risk systems embedded in regulated products from 2 August 2028. The official text, in
EUR-Lex, and the guidance of the national authorities, such as the Spanish AESIA, should be checked
for each new edition.
